\section{Translation into Algorithmic Research Language}

\subsection{General principle}

The project should not jump directly from discretionary language to machine learning. The correct sequence is:

\begin{enumerate}
\item event logging,
\item state reconstruction,
\item explicit feature definition,
\item simple rule-based detectors,
\item replay and live validation,
\item only then selective ML overlays.
\end{enumerate}

\subsection{Event families the repository must support}

\begin{longtable}{p{0.22\textwidth}p{0.25\textwidth}p{0.41\textwidth}}
\toprule
\textbf{Event family} & \textbf{Needed now} & \textbf{Why it matters} \\
\midrule
Order-book updates & yes & visible density, spread, queue geometry, depletion, refill, emptiness beyond level \\
Trade prints & yes & aggression, activity, executed pressure, absorption clues, ignition after break \\
Execution reports & later but schema-ready & live fill quality, slippage, queue realism, paper-versus-live comparison \\
Market status / session metadata & yes & open, close, halt, section changes, special regimes such as constrained price movement \\
Cross-instrument alignment & yes & leader-follower conditioning and day-context filters \\
\bottomrule
\end{longtable}

\subsection{Candidate state machines}

Most setups in the corpus translate more naturally into state machines than into static snapshots.

\paragraph{Bounce state machine.}
\begin{enumerate}
\item identify significant visible level,
\item observe approach,
\item measure aggressive attack intensity,
\item confirm hold or hidden replenishment,
\item enter on first validated rejection,
\item exit on first reliable obstacle or thesis failure.
\end{enumerate}

\paragraph{Breakout state machine.}
\begin{enumerate}
\item identify strong level,
\item detect repeated tests or compression,
\item confirm depletion of blocking liquidity,
\item confirm immediate post-break acceleration,
\item classify as continuation or failure within a short confirmation window.
\end{enumerate}

\paragraph{Robot state machine.}
\begin{enumerate}
\item detect repeated same-shape reposting,
\item estimate robot strength and freshness,
\item measure whether crowd is joining,
\item trade with the robot while the support remains active,
\item exit at first sign of robot stop or crowd overload.
\end{enumerate}

\subsection{Features suggested by the corpus}

\begin{itemize}
\item density size percentile within instrument-day,
\item distance from best bid/ask to nearest major density,
\item ratio of displayed size to executed volume at level,
\item refill count and refill cadence,
\item tape intensity in rolling short windows,
\item buy-versus-sell print imbalance,
\item forward book emptiness after a candidate breakout,
\item repeated test count of a level,
\item compression width before level attack,
\item realized travel after first touch,
\item immediate retrace after breakout,
\item first-hour performance of each setup family on the current day,
\item cross-instrument lead-lag metrics.
\end{itemize}

\subsection{What remains hard to encode}

Important parts of the corpus are still fuzzy:

\begin{itemize}
\item ``clean'' versus ``crooked'' feel,
\item recognition of when a level has already been overused,
\item subtle participant identity inference,
\item judgment that a move is ``already too late'',
\item understanding when a setup is crowded.
\end{itemize}

\researchnote{These are ideal candidates for later ML ranking or probability calibration \emph{after} explicit event sequences already exist. The first ML job should not be signal discovery from raw noise; it should be quality filtering for already-defined hypotheses.}

\subsection{Historical data limitations}

The repository's Russian historical dataset must be treated with caution:

\begin{itemize}
\item it is old,
\item it contains order book only,
\item it lacks prints and execution records,
\item many books are sparse or nearly empty,
\item reconstructed conclusions may be structurally misleading.
\end{itemize}

Therefore historical research on that data should focus on \emph{simple structural relationships}, for example:

\begin{itemize}
\item visible density versus short-horizon bounce probability,
\item density pull and immediate continuation,
\item repeated test count versus break probability,
\item compression near level versus post-break travel.
\end{itemize}

More nuanced discretionary claims about tape aggression, hidden liquidity, or crowding require fresh live capture.

\subsection{Live research roadmap implied by the corpus}

\begin{enumerate}
\item run continuous Bybit capture for public DOM and trades,
\item run T-Invest capture when the Russian market is open,
\item normalize and replay both feeds in the same internal event vocabulary,
\item annotate candidate setup events,
\item compare live observations against the formalized discretionary claims,
\item keep the first-hours adaptation statistics per instrument and per day,
\item only then consider ranking models or classifiers.
\end{enumerate}

\subsection{Repository architecture direction}

The proper architecture for this project is:

\begin{enumerate}
\item logging layer,
\item normalized event model,
\item replay engine,
\item feature layer,
\item rule-based setup detector layer,
\item analytics and journaling layer,
\item paper-trading layer,
\item live execution layer,
\item later ML overlay layer.
\end{enumerate}

This sequence preserves interpretability and makes it possible to tell whether a setup is failing because the idea is bad, because the implementation is wrong, or because the execution layer is distorting the result.

